\documentclass[11pt,a4paper]{article}

\usepackage{fontspec}
\usepackage[margin=22mm]{geometry}
\usepackage{amsmath,amssymb}
\usepackage{booktabs,longtable,array,tabularx}
\usepackage{graphicx}
\usepackage{caption}
\usepackage{xcolor}
\usepackage{hyperref}
\usepackage{enumitem}

\defaultfontfeatures{Ligatures=TeX,Scale=MatchLowercase}
\IfFontExistsTF{Noto Sans CJK KR}{
  \setmainfont{Noto Sans CJK KR}[
    ItalicFeatures={FakeSlant=0.18},
    BoldItalicFeatures={FakeSlant=0.18}
  ]
  \setsansfont{Noto Sans CJK KR}[
    ItalicFeatures={FakeSlant=0.18},
    BoldItalicFeatures={FakeSlant=0.18}
  ]
}{
  \IfFontExistsTF{Apple SD Gothic Neo}{
    \setmainfont{Apple SD Gothic Neo}[
      ItalicFeatures={FakeSlant=0.18},
      BoldItalicFeatures={FakeSlant=0.18}
    ]
    \setsansfont{Apple SD Gothic Neo}[
      ItalicFeatures={FakeSlant=0.18},
      BoldItalicFeatures={FakeSlant=0.18}
    ]
  }{
    \setmainfont{NanumGothic}[
      ItalicFeatures={FakeSlant=0.18},
      BoldItalicFeatures={FakeSlant=0.18}
    ]
    \setsansfont{NanumGothic}[
      ItalicFeatures={FakeSlant=0.18},
      BoldItalicFeatures={FakeSlant=0.18}
    ]
  }
}
\IfFontExistsTF{DejaVu Sans Mono}{
  \setmonofont{DejaVu Sans Mono}
}{
  \IfFontExistsTF{Menlo}{
    \setmonofont{Menlo}
  }{
    \setmonofont{Liberation Mono}
  }
}
\XeTeXlinebreaklocale "ko"
\XeTeXlinebreakskip = 0pt plus 1pt

\hypersetup{
  colorlinks=true,
  linkcolor=blue!50!black,
  urlcolor=blue!50!black,
  citecolor=blue!50!black,
  pdftitle={UUV MuJoCo v2.2 상세 설명 보고서},
  pdfauthor={OpenAI Codex}
}

\setlength{\parskip}{0.55em}
\setlength{\parindent}{0pt}
\setlength{\emergencystretch}{3em}
\renewcommand{\arraystretch}{1.22}
\newcolumntype{Y}{>{\raggedright\arraybackslash}X}
\newcommand{\code}[1]{\nolinkurl{#1}}
\captionsetup{
  font=small,
  labelfont=bf,
  justification=raggedright,
  singlelinecheck=false,
  skip=6pt
}

\title{UUV MuJoCo v2.2 상세 설명 보고서\\
\large 지금 코드만 보고도 구조와 물리를 설명할 수 있게 정리한 독립 문서}
\author{}
\date{2026-04-05}

\begin{document}
\maketitle
\tableofcontents
\clearpage

\section{이 문서의 목적}

이 문서는 \textbf{현재 repository의 active 코드 상태}를 기준으로, MuJoCo와 로봇/역학에 익숙하지 않은 사람도 바로 설명할 수 있도록 만든 \textbf{독립형 설명 문서}다. 기존 \code{mujoco\_fluid\_forces\_kr.tex} 를 수정한 것이 아니라, 2026년 4월 5일 시점의 \code{uuv\_mujoco/v2.2} 스냅샷을 별도 보고서로 다시 구성했다.

이 문서가 답하려는 질문은 단순하다.

\begin{enumerate}[leftmargin=2em]
  \item 지금 시스템은 어떤 파일들이 어떻게 연결되어 움직이는가.
  \item 로봇의 질량, 부력, 유체 저항, 스러스터, FCU 제어가 각각 어디에서 구현되는가.
  \item 이전 \code{custom} baseline과 현재 \code{ellipsoid} 기반 모델은 무엇이 다르고, 실제 측정 결과는 어떻게 다른가.
  \item 특정 이상 현상(느린 yaw, 깊이 흔들림, pitch bias, 과한 damping)이 보일 때 어디를 먼저 봐야 하는가.
\end{enumerate}

\section{한 장 요약}

\textbf{이 repository에는 두 runtime path가 남아 있지만, 설명의 기준 모델은 하나다.}

\begin{itemize}[leftmargin=2em]
  \item \code{custom}: 예전에 쓰던 Python 직접 계산형 6-DOF hydrodynamics baseline. 비교와 회귀 확인을 위해 남겨 둔 이전 모델이다.
  \item \code{ellipsoid}: \textbf{현재 모델}. scene의 \code{fluidshape="ellipsoid"} 와 \code{fluidcoef} 를 이용한 MuJoCo built-in ellipsoid fluid를 바탕으로, runtime에서 분산 질량/분산 부력/heave damping/reaction torque 같은 보정을 함께 얹는다.
\end{itemize}

즉 현재 조종감과 실제 비교의 기준은 \code{ellipsoid} 하나이고, \code{custom} 은 ``이전에는 이렇게 모델링했었다''를 설명하는 baseline으로 읽는 것이 맞다.

\textbf{MuJoCo는 강체 적분기이자 센서/scene 엔진이고, ArduSub는 제어기이며, JSON/XML은 물리 파라미터 파일이다.}

\begin{itemize}[leftmargin=2em]
  \item ArduSub는 ``얼마나 돌려라, 얼마나 내려가라'' 를 결정한다.
  \item MuJoCo runtime은 그 명령을 실제 힘으로 바꾼다.
  \item \code{sim\_profiles.json} 은 부력, drag, added mass, 관성 스케일을 정한다.
  \item \code{thruster\_params.json} 은 스러스터의 데드존, lag, 전역 gain, 반전 비대칭을 정한다.
  \item \code{tank\_ellipsoid\_scene.xml} 은 body, 센서 site, thruster site, built-in fluid proxy의 위치를 정한다.
\end{itemize}

\begin{figure}[ht]
  \centering
  \includegraphics[width=0.97\linewidth]{figures_v22/uuv_v22_architecture.png}
  \caption{현재 \code{v2.2} 전체 구조. 이 그림 하나만 보면 ``어떤 파일이 어떤 역할을 하는지''를 설명할 수 있다.}
\end{figure}

\section{물리엔진은 한 step에서 어떻게 모델링되는가}

이 문서에서 말하는 ``물리엔진''은 사실 두 층으로 나뉜다.

\begin{itemize}[leftmargin=2em]
  \item \textbf{Python runtime}: 어떤 힘과 모멘트를 만들어 MuJoCo에 넣을지 계산한다.
  \item \textbf{MuJoCo solver}: 그 힘과 모멘트, actuator, contact, built-in fluid를 받아 실제 강체 상태를 적분한다.
\end{itemize}

즉 현재 구현은 ``MuJoCo가 모든 유체력을 혼자 계산한다''도 아니고, ``Python이 모든 physics를 다 적분한다''도 아니다. \textbf{Python이 wrench를 만들고, MuJoCo가 적분한다}가 가장 정확한 설명이다.

\begin{figure}[ht]
  \centering
  \includegraphics[width=0.98\linewidth]{figures_v22/uuv_v22_runtime_step_pipeline.png}
  \caption{runtime 한 step의 실제 순서. Python runtime이 \code{data.ctrl} 와 \code{data.xfrc\_applied} 를 준비하고, MuJoCo가 \code{mj\_step} 으로 강체 적분과 built-in physics를 수행한다.}
\end{figure}

\subsection{실제로는 어떤 순서로 계산되는가}

한 step의 개념적 순서는 다음과 같다.

\begin{enumerate}[leftmargin=2em]
  \item MuJoCo state에서 현재 body pose, body/world velocity, depth, rotation matrix를 읽는다.
  \item FCU 또는 SITL servo 입력을 받아 deadzone, saturation, lag, force curve를 거쳐 \code{data.ctrl} 를 만든다.
  \item buoyancy point 또는 body component를 기준으로 잠김비를 계산하고, 분산 부력과 CoM/CoB torque를 만든다.
  \item 수면 근처 또는 완전 잠수 구간 heave damping을 더한다.
  \item 이전 baseline인 custom 경로라면 added mass, added-mass Coriolis, 선형 감쇠, 2차 감쇠로 body-frame hydrodynamic wrench를 계산한다.
  \item 현재 모델인 ellipsoid 경로라면 위 custom wrench는 끄고, 대신 MuJoCo built-in ellipsoid fluid가 scene의 proxy geom들에서 drag를 계산하게 둔다.
  \item thruster reaction torque, yaw torque scale 같은 보조 토크를 더한다.
  \item Python이 만든 wrench를 \code{data.xfrc\_applied} 에 넣고, actuator scalar는 \code{data.ctrl} 에 넣는다.
  \item 마지막에 MuJoCo가 \code{mj\_step(model, data)} 를 수행해 강체 적분, actuator 적용, built-in fluid/contact/constraint 해결을 한다.
\end{enumerate}

\subsection{좌표계는 왜 중요한가}

여기서 가장 중요한 구현 포인트는 \textbf{body frame과 world frame을 섞어 쓰지 않는 것}이다.

\begin{itemize}[leftmargin=2em]
  \item hydrodynamics helper가 쓰는 속도 상태는 \code{body frame}의 $\nu=[u,v,w,p,q,r]^T$ 이다.
  \item distributed buoyancy force는 계산 결과가 \code{world z} 방향 힘으로 들어간다.
  \item custom hydrodynamic wrench는 먼저 \code{body frame}에서 계산한 뒤, base rotation으로 world frame으로 돌려 \code{xfrc\_applied} 에 넣는다.
  \item MuJoCo의 \code{xfrc\_applied} 는 \code{[force\_xyz, torque\_xyz]} world frame 6-vector 이다.
\end{itemize}

즉 ``어떤 항을 어느 좌표계에서 계산하고, 마지막에 어디로 넣느냐''가 구현의 핵심이다. 이 부분을 놓치면 depth hold 중 수직 추력이 왜 수평 drift로 새는지 같은 현상을 설명할 수 없다.

\subsection{current model은 정확히 무엇이 MuJoCo이고 무엇이 Python인가}

현재 모델(\code{ellipsoid})을 기준으로 나누면 다음과 같다.

\begin{center}
\small
\begin{tabularx}{\linewidth}{>{\raggedright\arraybackslash}p{0.34\linewidth}YY}
\toprule
항목 & Python runtime이 담당 & MuJoCo가 담당 \\
\midrule
Distributed mass / inertia synthesis & \code{body\_components} 를 읽어 합성 CoM / inertia 계산 & 합성된 mass / inertia를 가진 강체 적분 \\
Thruster shaping & deadzone, command limit, lag, force scaling, yaw torque scale, reaction torque & actuator gear direction과 scalar 적용 \\
Buoyancy & submerged fraction, point-wise buoyancy, buoyancy torque 계산 & 주어진 force/torque로 rigid body 반응 \\
Custom hydrodynamics & added mass, Coriolis, linear/quadratic damping wrench 계산 & 없음 \\
Ellipsoid fluid drag & 없음 & scene의 \code{fluid\_*} proxy geom에서 built-in ellipsoid fluid 계산 \\
Rigid-body integration & 없음 & \code{mj\_step} 으로 pose/velocity/contact 적분 \\
Sensor state & publish용 변환 일부 & 실제 body 상태와 sensor frame 기반 원상태 제공 \\
\bottomrule
\end{tabularx}
\end{center}

\subsection{왜 이렇게 하이브리드 구조로 만들었는가}

이 구조를 택한 이유는 단순하다.

\begin{itemize}[leftmargin=2em]
  \item \textbf{MuJoCo만으로는} 분산 부력, custom added-mass Coriolis, runtime-specific yaw authority 같은 세밀한 수중 모델을 설명하기 어렵다.
  \item \textbf{Python만으로는} 안정적인 강체 적분, actuator, contact, sensor frame, visualization을 다시 만들 필요가 생긴다.
  \item 그래서 현재 구현은 ``수중 전용 force model은 Python에서 만들고, 최종 dynamics integration은 MuJoCo에 맡기는'' 형태로 정리되어 있다.
\end{itemize}

\section{용어를 모르면 먼저 여기부터}

\begin{center}
\small
\begin{tabularx}{\linewidth}{>{\raggedright\arraybackslash}p{0.18\linewidth}Y}
\toprule
용어 & 뜻 \\
\midrule
질량 (mass) & 병진 운동의 무거움을 정하는 값. 크면 앞뒤/좌우/상하로 움직이기 어려워진다. \\
관성 (inertia) & 회전 운동의 무거움을 정하는 값. 크면 roll/pitch/yaw가 시작도 느리고 멈추는 데도 오래 걸린다. \\
CoM & Center of Mass. 질량중심. 무게가 대표적으로 작용하는 점. \\
CoB & Center of Buoyancy. 부심. 부력이 대표적으로 작용하는 점. \\
부력 (buoyancy) & 잠긴 부피 때문에 생기는 위쪽 힘. 기본식은 $F_B = \rho g V_{\mathrm{sub}}$. \\
added mass & 물속에서 주변 유체도 같이 가속해야 해서 생기는 추가 관성. 수중체가 공기 중보다 더 무겁게 느껴지는 원인 중 하나다. \\
damping / drag & 이미 움직이는 물체를 줄여 주는 저항. 크면 명령을 끊었을 때 빨리 정리된다. \\
thruster gain & 조이스틱/FCU 명령이 실제 추력으로 얼마나 크게 바뀌는지 정하는 스케일. \\
lag ($\tau$) & 모터 명령이 즉시 힘으로 나오지 않고 지연되어 따라오는 정도. \\
fluidcoef & MuJoCo built-in ellipsoid fluid proxy에 주는 계수 묶음. drag와 angular drag 체감을 크게 바꾼다. \\
\bottomrule
\end{tabularx}
\end{center}

\section{현재 active 파일은 무엇이고, 각각 무엇을 바꾸는가}

\begin{center}
\small
\begin{tabularx}{\linewidth}{>{\raggedright\arraybackslash}p{0.34\linewidth}Y}
\toprule
파일 & 역할 \\
\midrule
\code{uuv\_mujoco/v2.2/launch\_uuv\_sim.sh} & MuJoCo runtime 실행 wrapper. scene 선택, fluid model 선택, ROS2/QGC bridge 옵션을 넘긴다. \\
\code{uuv\_mujoco/v2.2/reset\_uuv\_sim.sh} & MuJoCo, SITL, MAVProxy, QGC, 외부 MAVROS 잔존 프로세스를 정리한다. \\
\code{uuv\_mujoco/v2.2/start\_sitl\_mujoco\_mj311.sh} & reset 후 SITL과 MuJoCo를 같이 띄운다. 로그 디렉터리도 여기서 만든다. \\
\code{uuv\_mujoco/v2.2/start\_ardusub\_sitl\_mj311.sh} & ArduSub SITL 파라미터를 주입한다. yaw/depth/attitude PID와 RC deadzone도 여기서 바뀐다. \\
\code{uuv\_mujoco/v2.2/run\_urdf\_full.py} & 실제 runtime 핵심. thruster shaping, buoyancy, hydrodynamics, distributed mass/inertia, ROS2 bridge를 모두 처리한다. \\
\code{uuv\_mujoco/v2.2/config/sim\_profiles.json} & 이전 baseline(\code{custom})과 현재 모델(\code{ellipsoid})의 물리 knob 모음. 부력, drag, added mass, CoB offset, inertia scale이 여기 있다. \\
\code{uuv\_mujoco/v2.2/config/thruster\_params.json} & 스러스터 공통 deadzone, lag, global gain, 반전 비대칭, polynomial force curve가 여기 있다. \\
\code{uuv\_mujoco/v2.2/config/thruster\_performance.json} & T200 PWM-force curve 원본 데이터. 전압별 실제 측정 곡선이다. \\
\code{uuv\_mujoco/v2.2/scenes/tank\_custom\_scene.xml} & 이전 baseline인 custom hydrodynamics 경로용 scene. 비교/회귀 확인용으로 남아 있다. \\
\code{uuv\_mujoco/v2.2/scenes/tank\_ellipsoid\_scene.xml} & \textbf{현재 모델}의 핵심 scene. built-in ellipsoid fluid와 site, sensor, thruster, fluid proxy geom을 정의한다. \\
\code{uuv\_mujoco/v2.2/physics/sim\_profile\_helpers.py} & JSON profile을 runtime이 쓰기 쉬운 dataclass로 바꾸고, ellipsoid baseline을 계산한다. \\
\code{uuv\_mujoco/v2.2/physics/hydrodynamics\_helpers.py} & 잠김비, thruster shaping, polynomial force, 6-DOF hydrodynamics helper 함수가 있다. \\
\code{uuv\_mujoco/v2.2/physics/thruster\_mapping.py} & thruster 순서, ArduSub servo mapping, sensor FRD/FLU 기준을 정의한다. \\
\bottomrule
\end{tabularx}
\end{center}

\section{현재 로봇 형상과 좌표계는 어떻게 배치되어 있는가}

\subsection{장면의 핵심은 body 하나와 site 여러 개다}

현재 \code{tank\_ellipsoid\_scene.xml} 의 \code{base\_link} body 아래에는 세 부류의 위치 정보가 있다.

\begin{itemize}[leftmargin=2em]
  \item \textbf{thruster site}: 스러스터가 힘을 걸어 주는 위치
  \item \textbf{mass site}: 분산 질량 모델을 설명하기 위한 질량 component 위치
  \item \textbf{buoyancy site}: 분산 부력 모델을 설명하기 위한 부력 point 위치
\end{itemize}

여기서 중요한 점은, \textbf{scene XML의 mesh 자체가 질량을 계산하는 것이 아니라}, runtime이 JSON의 \code{body\_components} 를 읽어 \textbf{합성 CoM과 합성 inertia} 를 다시 계산한다는 것이다. 즉 mesh는 모양, JSON은 물리 파라미터, Python runtime은 실제 합성 계산기다.

\subsection{3D로 보면 무엇이 한 번에 보이는가}

\begin{figure}[ht]
  \centering
  \includegraphics[width=0.96\linewidth]{figures_v22/uuv_v22_vehicle_layout_3d.png}
  \caption{3D isometric layout. built-in ellipsoid fluid proxy geom, distributed mass component, buoyancy point, 센서 위치, thruster vector, body 축을 한 좌표계에 겹쳐 놓은 그림이다.}
\end{figure}

이 그림 하나로 바로 설명할 수 있는 것은 다음과 같다.

\begin{itemize}[leftmargin=2em]
  \item \textbf{질량은 어디에 분산되어 있는가}: 빨간 점 세 개가 중앙 인클로저와 좌우 하부 body의 질량 component다.
  \item \textbf{부력은 어디에 분산되는가}: 초록 점 네 개가 vertical thruster 바로 위의 buoyancy point다.
  \item \textbf{수직 추력은 어디서 올라오는가}: 검정 벡터가 vertical thruster plane을 보여 준다.
  \item \textbf{yaw 토크는 어떻게 만들어지는가}: 주황 벡터가 대각선 yaw thruster 방향을 보여 준다.
  \item \textbf{왜 trim bias가 생기는가}: CoM, 유효 buoyancy center, fluid proxy 중심, thruster plane의 상대 위치가 한 화면에서 보인다.
\end{itemize}

\subsection{위에서 보면 무엇이 보이는가}

\begin{figure}[ht]
  \centering
  \includegraphics[width=0.90\linewidth]{figures_v22/uuv_v22_vehicle_layout_top.png}
  \caption{상면도. 파란 타원은 built-in ellipsoid fluid proxy geom, 빨강은 질량 component, 초록은 부력 point, 검정 화살표는 yaw thruster 방향이다.}
\end{figure}

이 그림에서 바로 읽을 수 있는 핵심은 다음과 같다.

\begin{itemize}[leftmargin=2em]
  \item 로봇은 \textbf{중앙 인클로저 1개 + 좌우 하부 body 2개} 의 3분할 질량 모델을 쓴다.
  \item 부력은 중앙 한 점이 아니라 \textbf{수직 스러스터 바로 위 4점}에 나눠 준다.
  \item yaw thruster는 대각선 방향으로 놓여 있어, 순수 전진/순수 sway뿐 아니라 회전 모멘트도 만든다.
  \item 현재 \textbf{유효 buoyancy center}는 CoM의 x축과 거의 맞춘 상태다. 이건 depth hold 중 뒤로 밀리는 trim bias를 줄이기 위한 의도적인 정렬이다.
\end{itemize}

\subsection{옆에서 보면 무엇이 보이는가}

\begin{figure}[ht]
  \centering
  \includegraphics[width=0.92\linewidth]{figures_v22/uuv_v22_vehicle_layout_side.png}
  \caption{측면도. 질량, 부력, vertical thruster plane, 센서 위치, inertial reference를 함께 표시했다. pitch bias가 생기는 이유를 설명할 때 가장 직접적인 그림이다.}
\end{figure}

측면도에서 읽어야 할 것은 다음이다.

\begin{itemize}[leftmargin=2em]
  \item 깊이 유지 추력은 vertical thruster plane에서 나온다.
  \item 부력은 buoyancy point 또는 CoB offset을 통해 더 위쪽에서 작용한다.
  \item 질량중심 CoM과 부력중심 CoB, fluid drag center, thruster plane이 어긋나면 \textbf{depth hold가 pitch moment를 만든다.}
  \item 따라서 ``뒤로 미끄러진다''는 현상은 유체엔진 문제라기보다 \textbf{수직 추력이 body attitude를 통해 수평 성분으로 새는 trim 문제}인 경우가 많다.
\end{itemize}

\section{6DOF에서 어떤 물리력이 어떤 축에 들어가는가}

\subsection{먼저 body 축과 상태벡터를 읽어야 한다}

runtime이 쓰는 기본 상태벡터는

\begin{equation}
\nu = [u,\ v,\ w,\ p,\ q,\ r]^T
\end{equation}

이다. 여기서

\begin{itemize}[leftmargin=2em]
  \item $u,v,w$ 는 body frame 선속도(surge, sway, heave)
  \item $p,q,r$ 는 body frame 각속도(roll, pitch, yaw)
\end{itemize}

를 뜻한다. 즉 현재 simulator의 hydrodynamics는 결국 \textbf{6자유도 wrench}

\begin{equation}
\tau = [F_x,\ F_y,\ F_z,\ \tau_x,\ \tau_y,\ \tau_z]^T
\end{equation}

를 어떻게 만들 것인가의 문제로 볼 수 있다.

\begin{figure}[ht]
  \centering
  \includegraphics[width=0.95\linewidth]{figures_v22/uuv_v22_6dof_forces_3d.png}
  \caption{3D wrench view. simulator가 최종적으로 만들어 내는 여섯 개 채널, 즉 $F_x$, $F_y$, $F_z$, $\tau_x$, $\tau_y$, $\tau_z$ 를 입체적으로 표현한 그림이다.}
\end{figure}

\begin{figure}[ht]
  \centering
  \includegraphics[width=0.92\linewidth]{figures_v22/uuv_v22_6dof_axes.png}
  \caption{현재 runtime이 쓰는 body frame 6-DOF 상태 정의. hydrodynamics helper 함수와 added-mass Coriolis 계산은 모두 이 축 정의를 기준으로 한다.}
\end{figure}

즉 runtime의 모든 계산은 결국 다음 여섯 채널을 누가 얼마나 만드는가의 문제로 환원된다.

\begin{itemize}[leftmargin=2em]
  \item 선형 힘 3개: surge, sway, heave
  \item 회전 모멘트 3개: roll, pitch, yaw
\end{itemize}

\subsection{custom 경로에서는 어떤 힘을 합쳐서 최종 wrench를 만드는가}

현재 \code{run\_urdf\_full.py} 의 custom 경로는 개념적으로 다음 합을 만든다.

\begin{align}
\tau_{\mathrm{total}}
&=
\tau_{\mathrm{buoyancy}}
+ \tau_{\mathrm{surface\_heave}}
+ \tau_{\mathrm{added\ mass}}
+ \tau_{\mathrm{damping}}
+ \tau_{\mathrm{thruster\ reaction}}
\end{align}

여기서 더 핵심적인 custom 6-DOF hydrodynamic 항은

\begin{align}
\tau_{\mathrm{hydro}}
=
-M_A \dot{\nu}_r
-C_A(\nu_r)\nu_r
-D_1 \nu_r
-D_2 |\nu_r| \odot \nu_r
\end{align}

형태다.

\begin{itemize}[leftmargin=2em]
  \item $M_A$: added-mass diagonal
  \item $C_A(\nu_r)$: added-mass Coriolis matrix
  \item $D_1$: 선형 damping
  \item $D_2$: 2차 damping
  \item $\nu_r$: 상대 속도(body velocity에서 water current를 뺀 값)
\end{itemize}

즉 custom 경로는 ``drag를 하나 넣었다'' 수준이 아니라, \textbf{관성형 항, Coriolis형 항, 선형 감쇠, 2차 감쇠, hydrostatics}를 합쳐서 body wrench를 만든 뒤, 마지막에 world frame으로 돌려 \code{xfrc\_applied}에 더하는 구조다.

\begin{figure}[ht]
  \centering
  \includegraphics[width=0.97\linewidth]{figures_v22/uuv_v22_custom_wrench_pipeline.png}
  \caption{custom 6-DOF wrench pipeline. 실제 코드가 \code{data.xfrc\_applied} 에 쓰기 직전까지 어떤 항을 합치는지 개념적으로 정리한 그림이다.}
\end{figure}

\subsection{어떤 항이 어떤 자유도에 주로 작용하는가}

\begin{figure}[ht]
  \centering
  \includegraphics[width=0.98\linewidth]{figures_v22/uuv_v22_runtime_term_matrix.png}
  \caption{왼쪽은 각 물리 항이 어떤 DOF에 주로 작용하는지, 오른쪽은 custom / ellipsoid 경로 중 어느 경로에서 그 항이 실제로 켜지는지를 정리한 matrix다.}
\end{figure}

이 그림을 읽는 법은 다음과 같다.

\begin{itemize}[leftmargin=2em]
  \item \textbf{Distributed buoyancy force} 는 주로 heave($F_z$)에 들어간다.
  \item \textbf{Buoyancy torque from point offset} 은 roll/pitch torque에 주로 들어간다.
  \item \textbf{Surface heave damping} 은 수면 근처 z축 chatter를 줄이기 위한 항이다.
  \item \textbf{Added mass, added-mass Coriolis, linear damping, quadratic damping} 은 6개 축 모두에 들어갈 수 있다.
  \item \textbf{Yaw torque scale on yaw thrusters} 는 사실상 yaw 축 authority만 키우는 보조 항이다.
\end{itemize}

\subsection{이전 custom baseline과 현재 ellipsoid 모델은 무엇이 다르게 켜지는가}

현재 runtime에서 두 경로는 다음 차이를 가진다. 여기서 읽는 순서는 ``custom과 ellipsoid가 둘 다 현재 주력이다''가 아니라, \textbf{custom은 이전 baseline, ellipsoid는 현재 active 모델}이라는 전제를 두는 편이 맞다.

\begin{center}
\small
\begin{tabularx}{\linewidth}{>{\raggedright\arraybackslash}p{0.38\linewidth}ccY}
\toprule
항목 & custom & ellipsoid & 설명 \\
\midrule
distributed mass/inertia synthesis & on & on & 두 경로 모두 \code{body\_components} 로 합성 mass/CoM/inertia를 다시 쓴다. \\
distributed buoyancy force & on & on & 두 경로 모두 buoyancy point 또는 component point에서 부력을 나눠 적용한다. \\
buoyancy torque from offset & on & on & CoM과 buoyancy point 차이로 roll/pitch restoring torque를 만든다. \\
surface heave damping & on & on & 수면 근처 z축 chatter를 줄이기 위한 항이다. \\
full-submerged heave damping & off & on & 현재 ellipsoid 경로에만 완전 잠수 구간 heave damping 보정이 추가된다. \\
added-mass inertia / Coriolis & on & off & custom에서만 Python이 직접 계산한다. \\
linear / quadratic damping wrench & on & off & custom에서만 body-frame 6-DOF damping wrench를 직접 계산한다. \\
built-in ellipsoid fluid proxy force & off & on & ellipsoid 경로에서는 MuJoCo가 scene의 proxy geom들로 유체력을 계산한다. \\
thruster reaction torque & on & on & 스러스터 회전반작용 보정은 두 경로 모두 들어간다. \\
yaw torque scale & optional & on & 현재 설정에서는 ellipsoid만 \code{yaw\_torque\_scale=4.0} 를 쓴다. \\
\bottomrule
\end{tabularx}
\end{center}

즉 \textbf{현재 모델인 ellipsoid 경로는 ``scene만 믿는 순수 built-in mode''가 아니고}, MuJoCo built-in ellipsoid fluid 위에 runtime 보정 항을 덧붙인 \textbf{하이브리드 모델}이다. 반대로 \textbf{custom 경로는 현재 운영 모델이라기보다 이전 Python 직접 계산형 baseline}으로 보는 편이 맞고, rigid-body integration과 contact, sensor frame, visualization은 여전히 MuJoCo가 맡는다.

\section{질량, 관성, 부력은 현재 어떻게 구현되어 있는가}

\subsection{현재 질량 모델}

현재 두 profile 모두 질량 component는 동일하다.

\begin{center}
\small
\begin{tabular}{lccc}
\toprule
Component & mass [kg] & mass\_pos x [m] & mass\_pos z [m] \\
\midrule
center\_enclosure & 6.0 & -0.005 & 0.075 \\
port\_lower\_body & 4.5 & 0.000 & -0.005 \\
starboard\_lower\_body & 4.5 & 0.000 & -0.005 \\
\midrule
Total & 15.0 & \multicolumn{2}{c}{composite CoM = (-0.002, 0.000, 0.027)} \\
\bottomrule
\end{tabular}
\end{center}

즉 \textbf{중앙 상부에 6kg, 좌우 하부에 각각 4.5kg}이 있다고 생각하면 된다. 이 구조는 ``상부 인클로저 + 좌우 lower body'' 라는 실물 직관과 맞고, 따라서 초보자도 왜 roll/pitch 응답이 그런지 상상하기 쉽다.

\subsection{현재 관성 모델}

여기서 중요한 함정이 하나 있다. 현재 \code{tank\_ellipsoid\_scene.xml} 의 \code{diaginertia}는

\begin{verbatim}
diaginertia="0.001 0.001 0.001"
\end{verbatim}

로 매우 작다. 그러나 이 값이 최종 물리 모델의 전부는 아니다. 이 값은 \textbf{MuJoCo XML 로더가 scene을 받아들이기 위한 최소 안전값}이고, 실제 runtime에서는 \code{run\_urdf\_full.py} 가 \code{body\_components} 로 합성한 관성에 \code{body\_inertia\_scale\_xyz} 를 곱해 다시 덮어쓴다.

현재 계산 결과는 다음과 같다.

\begin{center}
\small
\begin{tabular}{lccc}
\toprule
구분 & $I_{xx}$ & $I_{yy}$ & $I_{zz}$ \\
\midrule
합성 raw inertia & 0.473376 & 0.370230 & 0.739326 \\
custom runtime inertia & 0.473376 & 0.370230 & 0.739326 \\
ellipsoid runtime inertia & 0.000473 & 0.000370 & 0.000739 \\
\bottomrule
\end{tabular}
\end{center}

왜 이렇게 되느냐면, 현재 \code{ellipsoid} profile 에서는

\begin{verbatim}
body_inertia_scale_xyz = [0.001, 0.001, 0.001]
\end{verbatim}

이기 때문이다. 즉 \textbf{회전 관성은 원래 합성값의 0.1\%만 남긴} 극단적 상태다. 이 값은 yaw를 매우 빠르게 만들 수 있지만, 동시에 damping과 controller가 과하면 ``관성은 거의 없는데 딱딱하게 붙잡는'' 느낌을 만들기도 한다.

\subsection{현재 부력 모델}

부력 기본식은 그대로다.

\begin{equation}
F_B = \rho g V_{\mathrm{sub}} \cdot \mathrm{buoyancy\_scale}
\end{equation}

하지만 이 문서에서 더 중요한 건 \textbf{어디에, 어떤 곡선으로, 어떤 비율로} 적용되느냐이다.

\begin{itemize}[leftmargin=2em]
  \item 잠김비 모델: \code{linear} 가 아니라 \code{ellipsoid}
  \item 부력 point: vertical thruster 위 4점
  \item point share: 각 0.25
  \item x/z offset: \code{cob\_x\_offset = 0.009}, \code{cob\_z\_offset = 0.010}
  \item buoyancy\_scale: \code{1.008}
\end{itemize}

즉 현재 유효 부력중심은 대략

\begin{equation}
(-0.002,\ 0.0001,\ 0.050)\ \mathrm{m}
\end{equation}

정도이며, composite CoM의

\begin{equation}
(-0.002,\ 0.0000,\ 0.027)\ \mathrm{m}
\end{equation}

보다 약 23mm 위에 있다. 따라서 roll/pitch에 대해서는 restoring moment가 생기지만, yaw에 대해서는 직접 restoring이 거의 없다.

\begin{figure}[ht]
  \centering
  \includegraphics[width=0.98\linewidth]{figures_v22/uuv_v22_buoyancy_model.png}
  \caption{현재 코드가 쓰는 잠김비/부력 모델. 선형 모델과 비교하면 ellipsoid 모델은 수면 중앙 근처에서 기울기가 더 크고, 양 끝에서는 더 부드럽게 포화된다.}
\end{figure}

이 그림이 의미하는 바는 명확하다.

\begin{itemize}[leftmargin=2em]
  \item 수면 근처에서는 작은 깊이 변화에도 부력 변화가 민감하다.
  \item 완전 잠수/완전 노출 근처에서는 기울기가 줄어들어 chatter를 줄인다.
  \item 따라서 z축 진동을 설명할 때 ``부력이 선형이냐 비선형이냐''는 의미가 있다.
\end{itemize}

\section{스러스터 모델은 현재 어떻게 구현되어 있는가}

\subsection{현재 공통 thruster knob}

\code{thruster\_params.json} 의 현재 공통 값은 다음과 같다.

\begin{center}
\small
\begin{tabular}{lc}
\toprule
항목 & 현재 값 \\
\midrule
deadzone & 0.002 \\
tau\_up & 0.04 \\
tau\_down & 0.06 \\
reverse\_asymmetry & 0.70 \\
reaction\_torque\_gain & 0.0012 \\
gain\_scale\_all & 2.6 \\
command\_limit & 0.90 \\
\bottomrule
\end{tabular}
\end{center}

\textbf{이 값들이 의미하는 바}는 다음과 같다.

\begin{itemize}[leftmargin=2em]
  \item \code{deadzone}: 스틱을 살짝 건드렸을 때 모터가 바로 울지 않게 하는 구간
  \item \code{tau\_up/down}: 힘이 갑자기 바뀌지 않고 1차 지연으로 따라오게 하는 시간상수
  \item \code{gain\_scale\_all}: 전체 추력을 한 번에 키우거나 줄이는 전역 배율
  \item \code{reverse\_asymmetry}: 후진 추력이 전진보다 약한 현실을 반영하는 계수
\end{itemize}

\subsection{실제 추력 생성 경로}

runtime에서 스러스터 명령은 다음 순서로 변환된다.

\begin{enumerate}[leftmargin=2em]
  \item pilot 또는 FCU 명령이 정규화된 \([-1, 1]\) command로 들어온다.
  \item \code{shape\_thruster\_command()} 가 deadzone과 command limit를 적용한다.
  \item \code{first\_order\_response()} 가 tau\_up/down에 따라 상태를 부드럽게 따라가게 만든다.
  \item measured T200 curve가 활성화돼 있으면 PWM-force curve를, 아니면 polynomial fallback force를 쓴다.
  \item 마지막에 per-thruster gain과 global gain을 곱해 실제 force가 된다.
\end{enumerate}

\begin{figure}[ht]
  \centering
  \includegraphics[width=0.98\linewidth]{figures_v22/uuv_v22_thruster_model.png}
  \caption{현재 스러스터 모델. measured T200 curve, deadzone/limit shaping, polynomial fallback force, first-order lag을 함께 보여 준다.}
\end{figure}

이 그림은 조종감을 설명할 때 가장 중요하다.

\begin{itemize}[leftmargin=2em]
  \item 저입력 구간이 너무 무디면 deadzone이나 polynomial 초반 기울기를 봐야 한다.
  \item 입력은 빠른데 실기보다 부드러우면 \code{tau\_up/down} 을 늘린다.
  \item 같은 제어기인데도 전체가 강하면 \code{gain\_scale\_all} 을 낮춘다.
\end{itemize}

\section{이전 custom baseline과 현재 ellipsoid 모델의 hydrodynamics는 어떻게 다른가}

\subsection{이전 custom baseline}

\code{custom} 경로는 \textbf{이전에 쓰던 기준 모델}로, Python runtime이 직접 계산한 hydrodynamic wrench를 \code{xfrc\_applied} 로 body에 주입한다. 여기에 들어가는 항은 크게 네 종류다.

\begin{itemize}[leftmargin=2em]
  \item hydrostatic buoyancy
  \item distributed buoyancy torque
  \item added mass / added-mass Coriolis
  \item linear / quadratic damping
\end{itemize}

\subsection{현재 ellipsoid 모델}

\code{ellipsoid} 경로는 \textbf{현재 모델}이다. scene에 있는 세 개의 fluid proxy geom이 MuJoCo built-in ellipsoid fluid를 계산하고, 그 위에 runtime에서 분산 질량, 분산 부력, heave damping, reaction torque, yaw torque scale 같은 보정이 더해진다. 현재 scene의 대표 값은 다음과 같다.

\begin{center}
\small
\begin{tabular}{lccc}
\toprule
Geom & position [m] & size [m] & fluidcoef \\
\midrule
fluid\_center\_enclosure & (-0.005, 0.000, 0.020) & (0.205, 0.115, 0.130) & (0.620, 0.175, 1.050, 0, 0) \\
fluid\_port\_lower\_body & (0.000, 0.208, 0.020) & (0.370, 0.080, 0.120) & (0.520, 0.150, 0.750, 0, 0) \\
fluid\_starboard\_lower\_body & (0.000, -0.208, 0.020) & (0.370, 0.080, 0.120) & (0.520, 0.150, 0.750, 0, 0) \\
\bottomrule
\end{tabular}
\end{center}

\subsection{profile에서 runtime으로 파생되는 hydrodynamic 값}

\begin{figure}[ht]
  \centering
  \includegraphics[width=0.98\linewidth]{figures_v22/uuv_v22_hydrodynamics_comparison.png}
  \caption{runtime이 실제로 쓰는 hydrodynamic 계수 요약. 이전 baseline인 \code{custom} 과 현재 모델인 \code{ellipsoid} 가 added mass와 damping을 어떻게 다르게 쓰는지 보여 준다.}
\end{figure}

현재 \code{ellipsoid} profile의 핵심 값은 아래처럼 읽으면 된다.

\begin{itemize}[leftmargin=2em]
  \item \code{effective\_cd\_linear = [0.100, 0.095, 0.145]}
  \item \code{effective\_cd\_angular = [2.70, 1.30, 0.005]}
  \item \code{added\_mass\_scale\_linear = [0.45, 0.28, 0.65]}
  \item \code{added\_mass\_scale\_angular = [1.00, 0.92, 1.02]}
  \item \code{linear\_damping\_ratio\_linear = 2.0}
  \item \code{linear\_damping\_ratio\_angular = 4.2}
\end{itemize}

여기서 특히 눈여겨볼 것은 \textbf{yaw angular drag의 유효 계수}와 \textbf{yaw torque scale} 이다. 현재 \code{effective\_cd\_angular[2]} 는 매우 낮고, 대신 \code{yaw\_torque\_scale = 4.0} 이라서 물리적으로는 yaw 저항을 낮추고, 추진 authority는 크게 키운 상태다.

\section{왜 MANUAL yaw와 DEPTH\_HOLD yaw 느낌이 다를 수 있는가}

초보자가 가장 헷갈리는 지점이 여기다. 같은 yaw stick인데 mode만 바꾸면 느낌이 달라지는 이유는, assisted mode에서는 yaw가 더 이상 direct motor command가 아니기 때문이다.

\begin{figure}[ht]
  \centering
  \includegraphics[width=0.98\linewidth]{figures_v22/uuv_v22_yaw_pipeline.png}
  \caption{MANUAL yaw와 assisted yaw의 차이. assisted mode에서는 yaw stick이 먼저 desired yaw rate로 바뀌고, 그다음 attitude controller를 거친다.}
\end{figure}

현재 injected SITL yaw 관련 파라미터는 다음과 같다.

\begin{center}
\small
\begin{tabular}{lc}
\toprule
파라미터 & 현재 값 \\
\midrule
RC4\_DZ & 10 \\
ACRO\_YAW\_P & 8.00 \\
ATC\_ANG\_YAW\_P & 2.20 \\
ATC\_ACC\_Y\_MAX & 2200.0 \\
ATC\_RATE\_Y\_MAX & 360.0 \\
ATC\_RAT\_YAW\_P/I/IMAX/D & 0.100 / 0.012 / 0.12 / 0.0025 \\
\bottomrule
\end{tabular}
\end{center}

즉 assisted mode yaw는 다음 네 층이 곱해진 결과다.

\begin{enumerate}[leftmargin=2em]
  \item RC deadzone과 stick scaling
  \item desired yaw rate 생성
  \item FCU yaw rate loop
  \item MuJoCo 쪽 actual yaw torque authority와 drag
\end{enumerate}

그래서 ``yaw가 느리다''고 느껴질 때는 \textbf{반드시 어느 층이 병목인지}를 분리해야 한다.

\section{현재 측정 결과는 무엇을 말해 주는가}

\subsection{실제 mode path 비교}

이 repository에는 이전 baseline인 custom과 현재 모델인 ellipsoid 두 경로에 대해 실제로 재생한 \code{ros2 bag} 요약 JSON이 이미 들어 있다. 여기서는 그 결과를 새 그림으로 다시 묶었다.

\begin{figure}[ht]
  \centering
  \includegraphics[width=0.97\linewidth]{figures_v22/uuv_v22_mode_paths_xy.png}
  \caption{mode별 measured XY path. custom과 ellipsoid가 같은 ``전진 $\rightarrow$ 90도 turn $\rightarrow$ 우측 이동'' 시퀀스를 어떻게 다르게 누적하는지 바로 볼 수 있다.}
\end{figure}

\begin{figure}[ht]
  \centering
  \includegraphics[width=0.97\linewidth]{figures_v22/uuv_v22_mode_metrics.png}
  \caption{path length, 최대 pitch disturbance, 90도 turn error 비교. 숫자만 보는 것보다 어떤 경로가 더 정리되는지 훨씬 직관적으로 드러난다.}
\end{figure}

이 결과가 의미하는 바는 다음과 같다.

\begin{itemize}[leftmargin=2em]
  \item 이전 baseline인 \code{custom}은 authority가 큰 대신 path가 더 길게 늘어지는 경향이 있다.
  \item 현재 모델인 \code{ellipsoid}는 이동량은 보수적이지만 turn error와 자세 disturbance가 더 작게 정리되는 경향이 있다.
  \item 따라서 보고서에서 ``현재 모델''이라고 부를 것은 ellipsoid 쪽이고, custom은 과거 모델 또는 비교용 baseline으로 설명하는 편이 맞다.
\end{itemize}

\subsection{저장돼 있는 archived engine snapshot}

기존 보고서에 들어가 있던 step-test 요약도 repository 안에 남아 있다. 이 값은 현재 mode-path 측정과는 별개의, \textbf{archived force/moment snapshot}이다.

\begin{figure}[ht]
  \centering
  \includegraphics[width=0.94\linewidth]{figures_v22/uuv_v22_engine_snapshot.png}
  \caption{repository에 저장돼 있는 archived engine snapshot. legacy custom, current custom, current ellipsoid의 step-test 지표를 세 축으로 다시 보여 준다.}
\end{figure}

이 그림은 딱 세 가지를 설명할 때 유용하다.

\begin{itemize}[leftmargin=2em]
  \item current custom은 가장 공격적인 이전 계열 모델이다.
  \item current ellipsoid는 지금 설명의 기준이 되는 현재 모델이다.
  \item legacy custom과 current custom은 모두 현재 모델 그 자체라기보다, ellipsoid로 오기 전 계열을 설명하는 비교 기준으로 읽는 편이 맞다.
\end{itemize}

\section{현재 문제를 볼 때 어디를 먼저 보면 되는가}

\begin{center}
\small
\begin{tabularx}{\linewidth}{>{\raggedright\arraybackslash}p{0.28\linewidth}>{\raggedright\arraybackslash}p{0.26\linewidth}Y}
\toprule
증상 & 먼저 볼 파일 & 먼저 볼 파라미터 \\
\midrule
yaw가 MANUAL보다 DEPTH\_HOLD에서 둔함 & \code{start\_ardusub\_sitl\_mj311.sh} & \code{RC4\_DZ}, \code{ACRO\_YAW\_P}, \code{ATC\_RAT\_YAW\_*} \\
yaw는 빠른데 금방 딱 멈춰 버림 & \code{sim\_profiles.json}, \code{start\_ardusub\_sitl\_mj311.sh} & \code{effective\_cd\_angular[2]}, \code{ATC\_ANG\_YAW\_P}, \code{ATC\_RAT\_YAW\_D} \\
depth hold 중 뒤로 밀림 & \code{sim\_profiles.json}, \code{tank\_ellipsoid\_scene.xml} & \code{cob\_x\_offset}, buoyancy point 위치, fluid proxy z 위치 \\
전진 시 과하게 숙여짐 & \code{tank\_ellipsoid\_scene.xml} & fluid proxy z, CoM/CoB 상대 위치, pitch damping \\
위아래로 hunting & \code{thruster\_params.json}, \code{start\_ardusub\_sitl\_mj311.sh}, \code{sim\_profiles.json} & \code{tau\_up/down}, \code{PSC\_D\_*}, \code{heave\_damping\_scale} \\
전체가 너무 둔함 & \code{thruster\_params.json}, \code{sim\_profiles.json} & \code{gain\_scale\_all}, \code{deadzone}, \code{added\_mass}, \code{effective\_cd} \\
병진은 좋은데 yaw만 약함 & \code{sim\_profiles.json}, \code{run\_urdf\_full.py} & \code{yaw\_torque\_scale}, yaw thruster gain, yaw drag \\
\bottomrule
\end{tabularx}
\end{center}

\section{실행과 재현 방법}

\subsection{권장 실행}

\begin{verbatim}
cd /Users/kanghyunmin/Desktop/uuv_sim/uuv_mujoco/v2.2
./reset_uuv_sim.sh --with-qgc-stop --wipe-eeprom
./start_sitl_mujoco_mj311.sh \
  --ros2-real-pkg-compat \
  --sitl-no-rebuild \
  -- \
  --tank-549x274x132 \
  --fluid-model ellipsoid \
  --profile ellipsoid \
  --no-qgc-video
\end{verbatim}

\subsection{새 문서 자산 생성}

\begin{verbatim}
cd /Users/kanghyunmin/Desktop/uuv_sim
/Users/kanghyunmin/.venvs/mujoco311/bin/python \
  document/docsource/generate_uuv_v22_explainer_assets.py
\end{verbatim}

\subsection{새 PDF 빌드}

\begin{verbatim}
cd /Users/kanghyunmin/Desktop/uuv_sim/document/docsource
xelatex -interaction=nonstopmode uuv_mujoco_v22_explainer_kr.tex
xelatex -interaction=nonstopmode uuv_mujoco_v22_explainer_kr.tex
\end{verbatim}

\section{이 문서만 보고 바로 설명할 때의 말하기 순서}

만약 발표나 회의에서 3분 안에 설명해야 한다면, 아래 순서를 그대로 쓰면 된다.

\begin{enumerate}[leftmargin=2em]
  \item ``이 시스템은 ArduSub가 제어기를 담당하고, MuJoCo runtime이 실제 힘과 센서를 담당하는 구조다.''
  \item ``물리는 \code{sim\_profiles.json}, \code{thruster\_params.json}, \code{tank\_ellipsoid\_scene.xml} 세 파일이 핵심이다.''
  \item ``질량은 3분할, 부력은 vertical thruster 위 4점 분산이다.''
  \item ``custom은 이전 baseline이고, 현재 모델은 built-in ellipsoid fluid에 runtime 보정을 얹은 ellipsoid 경로다.''
  \item ``실측 기준으로는 현재 ellipsoid 모델이 더 보수적이지만 더 정리되고, custom 계열은 이전 비교 기준으로 남아 있다.''
  \item ``문제가 생기면 먼저 증상별로 어느 층이 병목인지 분리해서 본다.''
\end{enumerate}

\section{결론}

현재 \code{uuv\_mujoco/v2.2} 는 더 이상 실험 파일이 뒤섞인 폴더가 아니라, \textbf{현재 모델(ellipsoid)과 이전 baseline(custom)을 함께 보존한 상태}다. 그리고 이 문서의 중요한 목적은 ``물리식을 많이 쓰는 것''보다, \textbf{지금 코드에서 어떤 파일이 어떤 현상을 만들고 있는지}를 한 번에 설명 가능하게 만드는 데 있다.

핵심만 다시 정리하면 이렇다.

\begin{itemize}[leftmargin=2em]
  \item scene XML은 위치와 shape를 준다.
  \item profile JSON은 부력, drag, added mass, 관성 스케일을 준다.
  \item thruster JSON은 입력을 힘으로 바꾸는 방식을 준다.
  \item SITL script는 FCU 제어 gain을 준다.
  \item runtime Python은 이 모든 것을 실제 wrench로 합성한다.
  \item 현재 모델은 \code{ellipsoid} 경로이며, MuJoCo built-in fluid 위에 runtime 보정 항을 더한 구조로 이해해야 한다.
  \item \code{custom} 은 현재 모델의 경쟁안이라기보다, 이전 계열 baseline으로 설명하는 것이 맞다.
\end{itemize}

따라서 누가 ``왜 이 로봇이 이렇게 움직이느냐''고 물으면, 이제는 감으로 답하지 않고 \textbf{scene, profile, thruster, FCU, runtime} 다섯 층으로 나눠 설명하면 된다. 이 문서는 바로 그 설명을 위해 만들어졌다.

\end{document}
